\documentclass[11pt]{article}
\usepackage{amsmath,amssymb,amsthm}
\newtheorem{lemma}{Lemma}
\newtheorem{theorem}{Theorem}
\title{Weak Faces and a Robust Signed Pair Obstruction}
\author{}\date{}
\begin{document}\maketitle
\paragraph{Status and scope.}
This is a complete written consequence of the saved exact rational
four-cone and quotient-bridge certificates. The new rank calculations
are checked by rational arithmetic, not Lean. It rules out a shared
signed residue pool for two specified targets on one pole hierarchy.
It does not rule out separately adaptive pools, adaptive pole grids,
the local compiler, or either global HeadComplexity target.

\begin{lemma}[Positive dependence with full rank]
Let $A(t)$ be a continuous real $m$ by $n$ matrix near $t=0$.
Suppose $A(0)$ has rank $n$ and there is a vector
$\lambda\in\mathbb R^m$ with every coordinate strictly positive and
$A(0)^\top\lambda=0$. Then for every sufficiently small $|t|$,
the only solution of $A(t)x\geq0$ is $x=0$.
\end{lemma}
\begin{proof}
Normalize $\lambda$ to have coordinate sum one. The augmented matrix
$M(0)=(A(0)^\top;\mathbf1^\top)$ has row rank $n+1$: the last row
cannot be a linear combination of the others, since multiplication by
$\lambda$ would give $1=0$. Choose $n+1$ columns making an invertible
square submatrix. Keep the other coordinates of $\lambda$ fixed and
solve the augmented equations on those columns. This gives a continuous
$\lambda(t)$ with $A(t)^\top\lambda(t)=0$, coordinate sum one, and
$\lambda(0)=\lambda$. It stays strictly positive for all small $t$.
The rank of $A(t)$ also stays $n$. If $A(t)x\geq0$, its positive
weighted sum is zero, so every entry is zero; full rank gives $x=0$.
\end{proof}

Fix $\epsilon=1/100$ and the sixteen lower poles
\[
 p_s=s_0\epsilon+s_1\epsilon^2+s_2\epsilon^4+s_3\epsilon^8,
 \qquad 0\leq s<16.
\]
Use the four sixteen-wire target words in
\texttt{butterfly-biased-cone-intersection.json}. Let $K_i$ be the
strict residue cone after eliminating its block biases, and let
$\overline K_i^{\rm weak}$ denote the cone obtained by replacing all
its strict comparison inequalities by weak inequalities.
Put $C_{st}=1/(1-p_s-p_t)$ and $v=C^{-1}\mathbf1$.

\begin{theorem}
The intersection of the four weak cones is $\{0\}$. For each of the
three pair partitions $(I,J)$, the weak quotient bridge
\[
 \{(a,\tau):a\in\bigcap_{i\in I}\overline K_i^{\rm weak},\quad
       a+\tau v\in\bigcap_{j\in J}\overline K_j^{\rm weak}\}
\]
is also $\{0\}$.
\end{theorem}
\begin{proof}
The saved four-cone dual uses seventeen comparison rows, all with
strictly positive weights. Their rank is sixteen. The saved dual for
each quotient bridge uses eighteen rows, all with strictly positive
weights, of rank seventeen in $(a,\tau)$. In each case a weakly
nonnegative vector of row evaluations must vanish, and full rank kills
the variable vector. These ranks and the exact positive dependences
are independently verified by
\texttt{butterfly-biased-weak-face-check.py}.
\end{proof}

\begin{theorem}[A robust signed parent pair]
Let $\Theta_0,\Theta_1$ be the two thirty-two-wire targets listed in
\texttt{butterfly-biased-cone-bridge.json}, under
\texttt{separating\_example.parent\_targets}. There is $\delta_0>0$
such that, for every $0<\delta<\delta_0$, these two targets do not
have a shared signed residue vector on the poles
\[
 p_s,\quad p_s+\delta\qquad(0\leq s<16),
\]
with query positions $1-p_s$ at every applicable node. Separate
target-dependent block biases are allowed. Residues may be unbounded
and may depend arbitrarily on $\delta$.
\end{theorem}
\begin{proof}
Take the first partition $I=\{0,1\}$, $J=\{2,3\}$.
The two parents have children $(0,2)$ and $(1,3)$ respectively,
and the same root word: eight minus signs followed by eight plus
signs in numerical prefix order. Write $\sigma_s$ for this root word.
At $\delta=0$, use child-indexed residue vectors $a,b\in\mathbb R^{16}$.
The root score vector is $C(a-b)$, and both parents have the same
root score. Thus a common root bias $\beta$ may be used whenever
their separate root biases are feasible.

The saved quotient dual supplies a vector $q$ with no zero coordinates,
$\sum_s q_s=0$, and $\sigma_s q_s<0$. Its weighted left child
comparison rows sum to $n=Cq$, and its weighted right rows sum to $-n$.
All eighteen weights are strictly positive. Lift these child rows to
the variables $(a,b,\beta)$, and append all sixteen signed root rows
\[
 \sigma_s\bigl(C_s,-C_s,1\bigr),\qquad 0\leq s<16.
\]
Give the root rows weights $|q_s|$. Their total is
$(-n,n,0)$, cancelling the child total $(n,-n,0)$ exactly.
All thirty-four weights are positive.

The lifted matrix has rank thirty-three. This follows either from
the checked rational rank or directly as follows. If all root rows
vanish, $a-b=-\beta v$. The vanishing child rows then form the
seventeen-variable quotient system, whose rank is seventeen by the
previous theorem. Hence $a=b=0$ and $\beta=0$.

For $\delta\ne0$, child-indexed residue coordinates assign $a_s$
to physical pole $p_s+\theta_s\delta$ and $b_s$ to the other pole,
where $\theta_s=(1-\sigma_s)/2$. This coordinate assignment is the
same for both parents because their root words agree. Lower child
comparisons and root inequalities are continuous rational functions
of $\delta$, with positive denominators for small $\delta$.
The preceding lemma applies to the thirty-four selected rows and
their thirty-three variables. Their weak feasible cone remains
$\{0\}$, so their strict feasible system is empty. Every proposed
shared pool for the full parent targets would satisfy this subsystem.
\end{proof}

\paragraph{What this does and does not close.}
The obstruction is stable at the fixed lower-pole scale
$\epsilon=1/100$ as the newest gap tends to zero. Thus a divergent
residue germ cannot rescue this particular common-pool induction,
and a normalized nonzero weak common face is absent. The theorem
does not assert the same behavior along the specific one-parameter
curve $\delta=\epsilon^{16}$ as $\epsilon\to0$. It also does not
assert that either parent target separately is infeasible. A universal
single-target decoder with independently chosen signed residues is
still unresolved.
\end{document}
